\begin{enumerate}
    \item[\textbf{Problem 5.1.}] If you were asked to conduct a population study, what would you be looking for, what would be known, and what would you want to know?

    \item[\textbf{Problem 5.2.}] What sort of assumptions would you make if you were asked to conduct a population study? On what basis could those assumptions be justified?

    \item[\textbf{Problem 5.3.}] What factors might restrain or otherwise influence the unrestrained growth seen in Figures 5.1 and 5.2?

    \item[\textbf{Problem 5.4.}] Determine the radius of the earth in both meters and feet from the surface areas given in Section 5.1. Are these values consistent with the conversion factors given in Table 2.3?

    \item[\textbf{Problem 5.5.}] Confirm that it would take almost two hundred thousand years to count a population of 5.63 million billion people at a rate of 1000 people/s.

    \item[\textbf{Problem 5.6.}] Confirm by differentiating eq. (5.4) that the $N(t)$ given therein satisfies eq. (5.1).

    \item[\textbf{Problem 5.7.}] How would the projections of Figures 5.1 and 5.2 change if the growth rate were, respectively, 1\% per year and 3\% per year?

    \item[\textbf{Problem 5.8.}] What annual growth rate would be needed to double one's money in seven years?

    \item[\textbf{Problem 5.9.}] Verify that eq. (5.12) is the solution to the exponential differential equation as given in eq. (5.10) by using the result that $\int du/u = \ln u + \text{constant}$.

    \item[\textbf{Problem 5.10.}] Show that the solution (5.12) to the differential equation (5.10) can also be found by assuming the following trial solution for $N(t)$: $N(t) = Ce^{\alpha t}$.

    \item[\textbf{Problem 5.11.}] Why is the proportionality constant in eq. (5.15) equivalent to having $\lambda < 0$ in eq. (5.10)? What sort of behavior would we then expect?

    \item[\textbf{Problem 5.12.}] Verify that each term in eq. (5.22) has the same physical dimensions.

    \item[\textbf{Problem 5.13.}] Confirm that eq. (5.23) is the correct solution for eq. (5.22).

    \item[\textbf{Problem 5.14.}] Use the definitions of resistance and capacitance to verify that the product $RC$ has the physical dimensions of time.

    \item[\textbf{Problem 5.15.}] Verify that eq. (5.29) was correctly derived from eq. (5.28).

    \item[\textbf{Problem 5.16.}] Apply Kirchhoff's Voltage Law (KVL) of eq. (5.28) to the circuit in Figure 5.7 to confirm the validity of eq. (5.27).

    \item[\textbf{Problem 5.17.}] Confirm by differentiating eq. (5.31) that it satisfies eq. (5.29).

    \item[\textbf{Problem 5.18.}] Calculate the initial slope of the charging capacitor from the solution given in eq. (5.31).

    \item[\textbf{Problem 5.19.}] Are eq. (5.35) and (5.38) related? How?

    \item[\textbf{Problem 5.20.}] Verify all of the steps that lead from eq. (5.35) to eq. (5.39).

    \item[\textbf{Problem 5.21.}] Construct a version of Table 5.3 for an annual interest rate of 18\%.

    \item[\textbf{Problem 5.22.}] If gasoline cost $0.70/gallon in 1978, calculate and compare the price inflation rates for the intervals 1933–1978 and 1978–2001.

    \item[\textbf{Problem 5.23.}] If the cost of money exceeds the cost of goods, what happens to $\lambda_{real}$?

    \item[\textbf{Problem 5.24.}] Speculate on the potential effects of $\lambda_{real}$ remaining negative for long periods of time.

    \item[\textbf{Problem 5.25.}] Look up the U.S. birth, death, immigration and emigration figures for the 10 decades of the 20th century and use the balance equation (5.45) to calculate the population changes that these rates predict.

    \item[\textbf{Problem 5.26.}] How do the predictions of Problem 5.25 compare with the actual U.S. census data?

    \item[\textbf{Problem 5.27.}] What are the implications for the model of eq. (5.48) of loosening the restriction that $\lambda_1, \lambda_2 > 0$?

    \item[\textbf{Problem 5.28.}] Confirm by differentiating eq. (5.52) that it satisfies eq. (5.51).

    \item[\textbf{Problem 5.29.}] Assuming that eq. (5.54b) can be solved for $F(t)$, show that $E(t)$ can be determined without further integration.

    \item[\textbf{Problem 5.30.}] Confirm that eq. (5.57) does follow from eq. (5.55) via eq. (5.56).

    \item[\textbf{Problem 5.31.}] Suppose you are given the solution for the enemy population that satisfies eq. (5.53) as $E(t) = E_0 \cosh \alpha t - \sqrt{\lambda_F / \lambda_E} F_0 \sinh \alpha t$, where $\alpha^2 = \lambda_E \lambda_F$. How long does it take for the enemy forces to be completely annihilated?

    \item[\textbf{Problem 5.32.}] Would the divide and conquer strategy work for a "linear attrition law" that, for equally effective armies, replaces eq. (5.57) with $F(t) - E(t) = F_0 - E_0$?

    \item[\textbf{Problem 5.33.}] Show that if it takes time $t_c$ to count a population $P(t)$ that has a growth rate of $\lambda$, the population will have increased by an amount equal to $\lambda t_c P(t)$.

    \item[\textbf{Problem 5.34.}] (a) If the population counting rate is $c$, how long does it take to count the population at time $t$? (b) How long does it take to count the population increase that occurred while it was being counted at time $t$?

    \item[\textbf{Problem 5.35.}] Find the actual world population numbers for 1970, 1980, 1990, and 2000. Use these data to update the projections shown in Figures 5.1 and 5.2.

    \item[\textbf{Problem 5.36.}] Using the 2\% growth rate, plot the world population from 1960 to 2060 with an ordinate scale of 10 billion people per 1.50 in. Does this curve look like "reasonable" exponential growth?

    \item[\textbf{Problem 5.37.}] The ordinate scales of Figures 5.1 and 5.2 are, respectively, 1.5 in = 100 billion people and 1.5 in = 3 million billion people. How much paper is needed to plot the 1960 world population of 3 billion people and the projected 2692 world population of 5.63 × $10^{15}$ people using each of these scales?

    \item[\textbf{Problem 5.38.}] Plot the growth of the world population from a 1960 value of 3 billion people at growth rates of 1\%, 2\%, and 3\% per year up to the year 2700 using semi-logarithmic paper. What shapes are these curves? What are their slopes and intercepts?

    \item[\textbf{Problem 5.39.}] What is the time constant, comparable to that of an RC circuit, for a population that decays at a fractional rate per unit time $\lambda$?

    \item[\textbf{Problem 5.40.}] How much must be set aside in 2002 in a savings account at 5.5\% per year to accumulate $1,000,000 by 2022? By 2042?

    \item[\textbf{Problem 5.41.}] Suppose there was a constant price inflation rate of 3\% per year, what would the investments of Problem 5.40 have to be to accumulate $1,000,000 in 2042 measured in 2002 dollars?

    \item[\textbf{Problem 5.42.}] The noted historian Stephen Ambrose chronicled the growth of the American railroads by listing the following amounts of total track per decade: 726 mi (1834), 4,311 (1844), 15,675 (1854), and 33,860 (1864). Determine: (a) the decade-by-decade growth rate; and (b) the exponential growth rate across all the data provided.

    \item[\textbf{Problem 5.43.}] Verify by differentiation and substitution that the following solution satisfies eqs. (5.53): $F(t) = F_0 \cosh \alpha t - \sqrt{\lambda_E / \lambda_F} E_0 \sinh \alpha t$ and $E(t) = E_0 \cosh \alpha t - \sqrt{\lambda_F / \lambda_E} F_0 \sinh \alpha t$.

    \item[\textbf{Problem 5.44.}] Confirm that the solution verified in Problem 5.43 satisfies Lanchester's square law of eq. (5.58).

    \item[\textbf{Problem 5.45.}] The initial strengths of two opposing armies are $F_0 = 10,000$ and $E_0 = 5,000$ troops, with equal loss rates of 0.1 per day. Who will win? How long will the battle take? (Hint: See Problem 5.31.) How many troops will the victor have when the enemy is vanquished? Graph the army populations until the enemy is completely annihilated.

    \item[\textbf{Problem 5.46.}] The initial strengths of two opposing armies are $F_0 = 10,000$ and $E_0 = 5,000$ troops, and $\lambda_F = 0.1$ per day. Who will win and with what remaining forces if $\lambda_E = 0.2$, $0.5$, and $1.0$ per day? What value of $\lambda_E$ would produce a draw?

    \item[\textbf{Problem 5.47.}] The landmark World War II battle of Iwo Jima began with troop sizes of $F_0 = 54,000$ and $E_0 = 21,500$ troops, with $\lambda_F = 0.0106$ per day and $\lambda_E = 0.0544$ per day. Absent any reinforcements, how long would this battle have lasted? How many troops would the victor have when the loser's forces were totally exhausted?

    \item[\textbf{Problem 5.48.}] To end the fighting for Iwo Jima in 28 days, how many troops would the United States have to have initially? How do the U.S. losses in this scenario compare to those found in the scenario of Problem 5.47?
\end{enumerate}
